\newpage
\section{Background} \label{sec:background}

Before diving into distributed training, let us briefly discuss the implementation and execution of local model training using PyTorch. Then, we explain and justify the idea of data parallelism and describe communication primitives.

\subsection{PyTorch}

PyTorch organizes values into \code{Tensor}s which are generic n-dimensional arrays with a rich set of data manipulating operations. A \code{Module} defines a transform from input values to output values, and its behavior during the forward pass is specified by its \code{forward} member function. A \code{Module} can contain \code{Tensor}s as parameters. For example, a \code{Linear} \code{Module} contains a \code{weight} parameter and a \code{bias} parameter, whose \code{forward} function generates the output by multiplying the input with the \code{weight} and adding the \code{bias}. An application composes its own \code{Module} by stitching together native \code{Module}s (\emph{e.g.}, linear, convolution, etc.) and \code{Function}s (\emph{e.g.}, relu, pool, etc.) in the custom \code{forward} function. A typical training iteration contains a forward pass to generate losses using inputs and labels, a backward pass to compute gradients for parameters, and an optimizer step to update parameters using gradients. More specifically, during the forward pass, PyTorch builds an autograd graph to record actions performed. Then, in the backward pass, it uses the autograd graph to conduct backpropagation to generate gradients. Finally, the optimizer applies the gradients to update parameters. The training process repeats these three steps until the model converges. 


\subsection{Data Parallelism}

PyTorch offers several tools to facilitate distributed training, including \code{DataParallel} for single-process multi-thread data parallel training using multiple GPUs on the same machine, \code{DistributedDataParallel} for multi-process data parallel training across GPUs and machines, and \code{RPC}~\cite{rpc} for general distributed model parallel training (\emph{e.g.}, parameter server~\cite{ps}). This paper focuses on \code{DistributedDataParallel}. Data parallelism enables distributed training by communicating gradients before the optimizer step to make sure that parameters of all model replicas are updated using exactly the same set of gradients, and hence model replicas can stay consistent across iterations. 

Parameter averaging is another popular technique to scale out model training. Similarly, it can launch multiple processes across multiple machines, but instead of synchronizing gradients, parameter averaging directly computes the average of all model parameters. This occurs after the local optimizer step, meaning that parameter averaging can be implemented completely as an auxiliary step and does not need to interact with local training steps at all, which is attractive as it can easily and cleanly decouple the code of distributed training and local iterations. There are several caveats with parameter averaging. 

\begin{itemize}
    \item Parameter averaging can produce vastly different results compared to local training, which, sometimes, can be detrimental to model accuracy. The root cause is that parameter averaging is \textbf{not} mathematically equivalent to processing all input data locally, especially when the optimizer relies on past local gradients values (\emph{e.g.}, momentum). As different model replicas are likely to see different gradients, the states in optimizers can gradually diverge, causing conflicting gradient descent directions. This can result in inexplicable differences in performance when switching from locally optimized models to large scale deployed models. 
    \item The structure of parameter averaging orchestrates computation (\emph{i.e.}, backward pass) and communication (\emph{i.e.}, computing average) into non-overlapping phases, using optimizer \code{step()} functions as a hard separation point. Regardless of how vigorously we optimize the computation or communication, one type of resource will stay idle at any given time instance, giving up a substantial performance optimization opportunity. 
\end{itemize}

Given the above fundamental pitfalls, we decided to implement distributed training using data parallelism to synchronize gradients instead of parameters. Note that, applications can still easily build parameter averaging using PyTorch. In fact, the collective communication feature described in Section~\ref{sec:collective} is an appropriate solution for this use case. Applications just need to explicitly launch \code{AllReduce} operations to calculate averaged parameters accordingly. 


\subsection{AllReduce}


\code{AllReduce} is the primitive communication API used by \code{DistributedDataParallel} to compute gradient summation across all processes. It is supported by multiple communication libraries, including NCCL~\cite{nccl}, Gloo~\cite{gloo}, and MPI~\cite{mpi}. The \code{AllReduce} operation expects each participating process to provide an equally-sized tensor, collectively applies a given arithmetic operation (\emph{e.g.}, \code{sum}, \code{prod}, \code{min}, \code{max}) to input tensors from all processes, and returns the same result tensor to each participant. A naive implementation could simply let every process broadcast its input tensor to all peers and then apply the arithmetic operation independently. However, as \code{AllReduce} has significant impact on distributed training speed, communication libraries have implemented more sophisticated and more efficient algorithms, such as ring-based \code{AllReduce}~\cite{nccl} and tree-based \code{AllReduce}~\cite{nccl:tree_reduce}. As one \code{AllReduce} operation cannot start until all processes join, it is considered to be a synchronized communication, as opposed to the P2P communication used in parameter servers~\cite{ps}.